\part{Picard 概形：一般性质}\label{fga3.vi}\thispagestyle{fancy}\begin{original}{}\label{fga3.vi-original-citation}\skipsinglepar\par{}A. Grothendieck．“Technique de descente et théorèmes d'existence en géométrie algébrique, VI: Les schémas de Picard: Propriétés générales”．\emph{Séminaire Bourbaki} \textbf{14}（1961–62），第 236 次报告，1962 年 5 月。\url{http://www.numdam.org/book-part/SB_1961-1962__7__221_0/}\end{original}\setcounter{chapter}{-1}\chapter[{对上次报告（[3]，V）的补充}]{对上次报告（[3]，V）的补充\textsuperscript{(1)}}\thispagestyle{fancy}\label{fga3.vi-0}\par{}\emph{［译校注］原刊标题保留了脚注标记(1)，但本库法文汇编扫描及 Numdam 原刊扫描中均未定位到相应脚注正文。此处保留标记，未以推测文字补造原注。}\par{}对于 \hyperref[fga3.v]{[3]，V} 中提出的 Picard 预概形的存在性问题，已有一些进展：\begin{enumerate}
    
 \item[a.]（Mumford。）若 \(f\colon  X\to  S\) 是射影可分态射（即平坦且纤维可分），则 Picard 预概形 \(\underline {\operatorname {Pic}}_{X/S}\) 存在，这在一般情形中并不成立；即使 \(f\) 的纤维维数为 \(1\)，且 \(S\) 是完备离散赋值环的谱，也是如此。一个反例如下：取 \(S=\operatorname {Spec}\mathbb {R}[{[t]}]\)，取 \(X\) 为 \(\mathbb {P}_S^2\) 中由方程 \(x^2+y^2=tz^2\) 定义的子概形（齐次变量为 \(x,y,z\)）。它表示一条退化为几何上相交的两条直线的圆锥曲线，但在域 \(\mathbb {R}\) 上的特殊纤维却不可约（它在 \(\mathbb {R}\) 中由方程 \(x^2+y^2=0\) 给出）。容易看出，作 étale 扩张 \(S'\to  S\)（其中 \(S'=\operatorname {Spec}\mathbb {C}[{[t]}]\)）后，\(X'/S'\) 的 Picard 预概形存在，并由此将其明确描述为若干 \(\widetilde {S}\) 的副本之和，其中 \(\widetilde {S}\) 是将 \(S\) 的原点复制无穷多次所得。容易看到，\(\underline {\operatorname {Pic}}_{X'/S'}\) 上相对于 \(S'\to  S\) 的下降数据（这里由 \(S'\) 在 \(S\) 上的 Galois 群 \(\mathbb {Z}/2\mathbb {Z}\) 的作用给出）并不有效，因为这个群置换某些被复制的点（因而有些轨道不包含在任何仿射开集中）。不过，Mumford 已证明：若 \(f\colon  X\to  S\) 是可分射影态射，且对所有 \(s\in  S\)，\(X_s\) 的不可约分支相对于 \(k(s)\) 都几何不可约，则 \(\underline {\operatorname {Pic}}_{X/S}\) 存在。证明基于他关于取商定理的一个精细化；参见 \hyperref[fga3.vi-bib-7]{[7]}。还要注意，即使不对纤维 \(X_s\) 的不可约分支作任何假设，下面将引入的概形 \(\underline {\operatorname {Pic}}_{X/S}^\tau \) 仍有可能总是存在。


    
 \item[b.]（Murre。）若 \(X\) 是一个域上的固有概形，则 \(\underline {\operatorname {Pic}}_{X/k}\) 存在。证明部分地利用了 Chevalley 的证明 \hyperref[Che1960]{[1]}，而根本上依靠 Picard 函子的群结构。

  \end{enumerate}关于 Picard 概形理论的其他一些评注，特别是它与 Abel 概形的关系，建议参阅 \hyperref[fga3.vi-bib-7]{[7]}。最后，本报告一个明显的不足是缺少“等价判据”，即能够比较射影概形及其超平面截面的 Picard 概形的判据。建立这些判据所需的关键定理可见 \hyperref[fga3.vi-bib-5]{[5]}，并须将它们与 Picard 概形的存在定理结合。\chapter{交换群预概形的拓扑性质}\thispagestyle{fancy}\label{fga3.vi-1}\par{}设 \(k\) 是域，\(G\) 是 \(k\) 上的群预概形。单位元 \(e\) 是 \(k\)-有理点，因而必为闭点，立刻可得 \(G\times _k G\) 的对角线闭，从而 \(G\) 是 \emph{分离} 的：\emph{域上的每个群预概形都是分离的}。用 \(G^0\) 表示单位元 \(e\) 的连通分支。因为 \(e\) 是 \(k\)-有理点，\(G^0\) 实际上是 \emph{几何连通} 的，即 \(G\mapsto  G^0\) 与到另一个域的基变换相容。还可得出，\(G^0\) 在乘法下稳定（集合意义下）；若 \(G\) 局部 Noether，则 \(G^0\) 为开集，于是可将 \(G^0\) 视为 \(G\) 的一个 \emph{开子群}。以下假设 \(G\) 在 \(k\) 上局部有限型；则 \(G^0\) 是 \emph{几何不可约且在 \(k\) 上有限型} 的。事实上，可以假设 \(k\) 代数闭，进而假设 \(G\) 约化（此时 \(G_\mathrm {red}\) 是 \(G\) 的子群，因为 \(G_\mathrm {red}\times _k G_\mathrm {red}\) 此时也约化）。于是它在某个非空开集上在 \(k\) 上光滑；平移此开集，便知它处处光滑。但这样一来，\(G\) 局部不可约，所以其不可约分支就是连通分支，因而 \(G^0\) 不可约。现在设 \(U\) 是 \(G^0\) 中 \(e\) 的一个仿射邻域；利用 \(G^0\) 不可约，立刻可知 \(U\cdot  U=G^0\)，这就证明 \(G^0\) 拟紧，因而在 \(k\) 上有限型。\par{}为简单起见，假设 \(G\) 交换。对每个整数 \(n>0\)，以 \(G^{(n)}\) 表示 \(G^0\) 在\(G\) 中的 \(n\) 次幂同态 \(\varphi _n\)下的逆像，于是 \(G^{(n)}\) 是 \(G\) 的开子群。令 \[     \begin {aligned}       G^\tau        &= \bigcup _{n>0} G^{(n)}     \\G^\sigma        &= \bigcup _{(n,p)=1} G^{(n)}     \\G^\rho        &= \bigcup _{h>0} G^{(p^h)}     \end {aligned}   \]，其中 \(p\) 是域 \(k\) 的特征指数。这些都是 \(G\) 的开子群，且满足 \[     \begin {aligned}       G^\sigma  \cap  G^\rho        &= G^0     \\G^\sigma  \cdot  G^\rho        &= G^\tau .     \end {aligned}   \]。\begin{remark}{}\label{fga3.vi-1-remark}\skipsinglepar\par{}可以构造商群概形 \(G/G^0=N\)（参见 \hyperref[fga3.iv]{[3]，IV}），并将 \(G^\tau \)、\(G^\sigma \)、\(G^\rho \) 分别定义为 \(N\) 的挠子群（分别地，其 \(p\)-准素分量，以及其 \(p\)-准素分量的自然补群，即对满足 \(q\neq  p\) 的素数 \(q\) 取各 \(q\)-准素分量之和）在 \(G\) 中的逆像。注意，\(N\) 是 \(k\) 上可分的离散群概形，因此（一旦选定 \(k\) 的代数闭包 \(\overline {k}\)，从而得到 Galois 群 \(\pi \)）可将它等同于一个通常的离散群，其上有 \(\pi \) 通过自同构的作用。挠子群的构造及其分解为各 \(q\)-准素分量，可以用这一方式作显然的解释。若 \(G\) 是 \(k\) 上固有概形 \(X\) 的 Picard 概形，则可以将 \(N\) 称为 \(X\) 在 \(k\) 上的（约化）\emph{Néron–Severi 概形}。若 \(G_\mathrm {red}^0\) 是 \(G^0\) 的群子概形——特别地，当 \(k\) 为完美域，或 \(G^0\) 在 \(k\) 上固有时（例如 \(X\) 几何正规时）即是如此——则还可以引入商 \(N'=G/G_\mathrm {red}^0\)。从特殊化的观点，即当 \(X\) 在代数概形族中变化时看，它的表现往往比 \(N\) 更好。\end{remark}\par{}现在设 \(S\) 是局部 Noether 预概形，\(G\) 是 \(S\) 上的群预概形，且在 \(S\) 上局部有限型。不假设 \(G\) 在 \(S\) 上有限型，也不假设它在 \(S\) 上分离。令 \[     G^0     = \bigcup _{s\in  S} (G_s)^0   \]；若 \(G\) 交换，再令 \[     \begin {aligned}       G^\tau        &= \bigcup _{s\in  S} (G_s)^\tau      \\G^\sigma        &= \bigcup _{s\in  S} (G_s)^\sigma      \\G^\rho        &= \bigcup _{s\in  S} (G_s)^\rho .     \end {aligned}   \]。这些是 \(G\) 的子集，在 \(G\) 的乘法下稳定，但这并不显然蕴含它们可以由 \(G\) 的子群预概形定义。特别地，似乎 \emph{一般并不存在底层集合为 \(G^0\) 的 \(G\) 的子群预概形}。当然，若其中某个集合是开集，赋予诱导结构后，它就是 \(G\) 的开子群预概形。我们将看到，对于 \(G^\tau \)，情形总是如此；因此，从可表函子的观点看，特别有 \emph{从“特殊化”的观点看，数值等价比代数等价表现得更令人满意}。以上所定义的集合的主要一般性质如下：\begin{theorem}{1.1}\setnumber{1.1}\label{fga3.vi-1-theorem-1.1}\skipsinglepar\par{}\(G^0\)、\(G^\tau \)、\(G^\sigma \) 和 \(G^\rho \) 都是 \emph{局部可构造} 的。此外：\begin{enumerate}
      
 \item[i.]\(G^0\) 在 \(S\) 上拟紧。若各 \(G_s^0\) 固有，且 \(G\) 在 \(S\) 上分离，则 \(G^0\) 在 \(S\) 上固有，因而在 \(G\) 中闭。

      
 \item[ii.]\(G^\tau \) 为开集。若 \(G^0\) 为闭集，则 \(G^\tau \) 亦然。

      
 \item[iii.]若 \(G^0\) 为闭集，则只要处于等特征情形，即 \(S\) 的所有剩余域有相同特征，\(G^\sigma \) 也为闭集。若 \(G^0\) 为闭集，且 \(G\to  S\) 在 \(G^\sigma \) 的各点处普遍开（参见下面的 \hyperref[fga3.vi-1-corollary-1.5]{推论~\ref*{fga3.vi-1-corollary-1.5}}），则 \(G^\sigma \to  S\) 普遍开。

      
 \item[iv.]若 \(G^0\) 为闭集，则 \(G^\rho \) 亦然。若处于等特征情形，且对每个满足 \((n,p)=1\) 的整数 \(n>0\)，\(G\) 中的 \(n\) 次幂同态为开映射，则 \(G^\rho \) 为开集。

    \end{enumerate}\end{theorem}\par{}现在给出证明的一些提示。\(G^0\) 局部可构造且在 \(S\) 上拟紧，这一事实包含在以下引理中：\begin{lemma}{1.2}\setnumber{1.2}\label{fga3.vi-1-lemma-1.2}\skipsinglepar\par{}若 \(S\neq \varnothing \)，则 \(S\) 中存在非空开集 \(U\)、\(U\) 上纤维连通的有限型群概形 \(H\)，以及一个群概形同态 \(H\to  G|U\)，后者是像为 \(G^0|U\) 的开浸入。\begin{proof}\par{}为证明此引理，可假设 \(S\) 不可约；设 \(\eta \) 为其泛点。作基变换 \(S'=\operatorname {Spec}(\mathcal {O}_{S,\eta })\to  S\)，便得到局部 Artin 环 \(\mathcal {O}_{S,s}=A\) 上的群概形 \(G'\)，其中有一个 \emph{开} 群子概形 \({G'}^0\)，它 \emph{在 \(A\) 上有限型}，如上所述。因此，它来自 \(\eta \) 的某个开邻域 \(U\) 上的有限型群概形 \(H\)，而典范浸入 \({G'}^0\to  G'\) 来自一个开浸入 \(H\to  G|U\)；当 \(U\) 足够小时，后者是群概形同态。若 \(U\) 足够小，\(H\) 的纤维连通；若 \(U\) 足够小，这些纤维均有相同维数，即 \(G\) 的纤维的维数。因此，当 \(U\) 足够小时，对所有 \(s\in  U\)，\(H_s\) 在 \(G_s\) 中的像恰为 \(G_s^0\)。引理得证。\end{proof}\end{lemma}\par{}\hyperref[fga3.vi-1-theorem-1.1]{定理~\ref*{fga3.vi-1-theorem-1.1}} 的 (i) 中的第二个断言包含在下列引理中（将此引理用于 \(G\) 中 \(G^0\) 的一个拟紧开邻域）：\begin{lemma}{1.3}\setnumber{1.3}\label{fga3.vi-1-lemma-1.3}\skipsinglepar\par{}设 \(X\) 是 \(S\) 上的有限型分离预概形，\(S\) 局部 Noether；设 \(g\) 是 \(X\) 在 \(S\) 上的截面，\(X^0\) 是各 \(g(s)\) 在各 \(X_s\) 中的连通分支之并。设 \(s\in  S\) 使得 \(X_s^0\) 在 \(k(s)\) 上固有。那么，\(s\) 有一个开邻域 \(U\)，使 \(X^0|U\) 在 \(U\) 上固有，从而也在 \(X|U\) 中闭。\begin{proof}\par{}对底作忠实平坦下降，可归约到 \(S\) 是完备局部环的谱、\(s\) 为其闭点的情形。应用 [\hyperref[GD1960]{2}，III，§5.5.1]，可见 \(X\) 分解为两个不交开集 \(X'\) 和 \(X''\) 的和，其中 \(X'\) 在 \(S\) 上固有，并且 \(X'_s=X_s^0\)。这使我们可归约到 \(X=X'\) 的情形，即 \(X\) 在 \(S\) 上固有的情形。此时可运用标准证明，利用子集固有性的赋值判据（[\hyperref[GD1960]{2}，II] 中遗漏了这一判据）。\end{proof}\end{lemma}\par{}现在证明 \(G^\tau \) 是开集；考虑到 \(G^\tau \) 的构造（与 \(G^0\)、\(G^\sigma \)、\(G^\rho \) 的构造一样）与基扩张交换，这等价于证明：对于 \(G\) 在 \(S\) 上的每个截面 \(g\)，\(g^{-1}(G^\tau )\) 是开集。这要求证明两点：\begin{enumerate}
    
 \item[a.]设 \(y\in  S\) 满足 \(g(y)\in  G^\tau \)。那么，对 \(y\) 的所有邻近点 \(y'\in \overline {y}\)，都有 \(g(y')\in  G^\tau \)。

    
 \item[b.]设 \(y'\in  S\) 满足 \(g(y')\in  G^\tau \)。那么，对 \(y'\) 的每个泛化 \(y\)，都有 \(g(y)\in  G^\tau \)。

  \end{enumerate}对于 (a)，注意存在整数 \(n>0\) 使 \(g^n(y)\in  G^0\)；由于 \(G^0\) 可构造，\((g^n)^{-1}G^0\) 也可构造；于是对 \(y\) 的所有邻近点 \(y'\in \overline {y}\)，都有 \(g^n(y')\in  G^0\)。对于 (b)，注意各 \(g(y')^n=g^n(y')\) 都留在 \(G_{y'}\) 的某个拟紧开集中（因为它们包含在模 \(G_{y'}^0\) 的有限多个类中）；因此 \(G\) 中存在一个拟紧开集 \(U\)，包含各 \(g^n(y')\)，从而也包含其泛化 \(g^n(y)\)。所以，\(g(y)\) 的各次幂都留在 \(G_y\) 的某个拟紧开集中，容易由此推出 \(g(y)\in  G_y^\tau \)。\par{}假设 \(G^0\) 为闭集；我们将证明 \(G^\tau \) 也是闭集。已知 \(G^\tau \) 为开集，因而局部可构造，所以只须证明它在特殊化下稳定。这是因为它是一族闭子集的并，这些闭子集就是闭子集 \(G^0\) 在各 \(n\) 次幂同态 \(\varphi _n\) 下的逆像。\par{}同一论证表明，只要 \(G^0\) 为闭集，\(G^\sigma \) 和 \(G^\rho \) 也都是闭集（对前者还须附加等特征假设）；只需先证明 \(G^\sigma \) 和 \(G^\rho \) 局部可构造。但对于 \(x\in  G^\tau \)，令 \(\nu (x)\) 为使 \(n\) 次幂同态 \(\varphi _n\) 将 \(x\) 映入 \(G\) 的最小整数 \(n>0\)。则 \(G^\sigma \)（分别地，\(G^\rho \)）由满足 \(\nu (x)\) 与 \(p\) 互素（分别地，为 \(p\) 的某个幂）的那些 \(x\in  X\) 组成。因而，所断言的可构造性由下面更精确的引理得到：\begin{lemma}{1.4}\setnumber{1.4}\label{fga3.vi-1-lemma-1.4}\skipsinglepar\par{}\(G^\tau \) 上的函数 \(\nu \) 局部可构造。\begin{proof}\par{}这就是说，对每个整数 \(n>0\)，满足 \(\nu (x)=n\) 的 \(x\in  G^\tau \) 所成的集合局部可构造。但这个集合是 \(\varphi _n^{-1}(G^0)\) 与各 \(\varphi _d^{-1}(G^0)\) 之并的差集，其中 \(d\) 遍历 \(n\) 的真因子。由于 \(G^0\) 局部可构造，各 \(\varphi _i^{-1}(G^0)\) 也都局部可构造，故上述差集亦然。\end{proof}\end{lemma}\par{}假设 \(G^0\) 闭，且 \(G\to  S\) 在 \(G^\sigma \) 的各点处普遍开；我们将证明 \(G^\sigma \to  S\) 为开映射，即将任意 \(x\in  G^\sigma \) 在 \(G^\sigma \) 中的邻域映为 \(y=f(x)\) 的邻域。由于 \(G^\sigma \) 局部可构造，只须证明：对 \(y\) 的每个泛化 \(y'\)，在 \(y'\) 上方存在 \(x\) 在 \(G\) 中的一个泛化 \(x\)'。作基变换，可归约到 \(S\) 是离散赋值环的谱，\(y\) 和 \(y'\) 分别为闭点和泛点的情形。利用 \(G\to  S\) 在 \(x\) 处为开映射（因而在 \(y'\) 上方存在 \(x\) 在 \(G\) 中的泛化 \(x_1\)），还可以（再作一次基变换后）假设存在 \(G\) 在 \(S\) 上的截面 \(g\)，使 \(x=g(y)\)。若 \(k(y')\) 的特征为 \(0\)，只须任取 \(x\) 在 \(G\) 中位于 \(y'\) 上方的泛化 \(x'\)；由于 \(G^\tau \) 为开集，它属于 \(G^\tau \)，进而由 \(G_{y'}^\sigma =G_{y'}^\tau \) 可知它属于 \(G^\sigma \)。若 \(k(y')\) 的特征为 \(p>0\)，则令 \[     \nu (g(y'))     = p^hm     \qquad \text {其中 }(m,p)=1   \]，取满足 \(ap^h+bm=1\) 的整数 \(a\) 和 \(b\)，并令 \[     \begin {aligned}       g_1       &= g^{ap^h}     \\g_2       &= g^{bm}     \end {aligned}   \]，于是 \(g=g_1g_2\)。由构造，有 \(g_1(y')\in  G^\sigma \) 和 \(g_2(y')\in  G^\rho \)。由于 \(G^0\) 闭，推出 \(g_1(y)\in  G^\sigma \) 和 \(g_2(y)\in  G^\rho \)；又因 \[     g(y)     = g_1(y)g_2(y) \in  G^\sigma    \]，还有 \(g_2(y)\in  G^\sigma \)，所以 \[     g_2(y)\in  G_y^\sigma \cap  G_y^\rho      = G_y^0.   \]。然而，由假设及 \(S\) 为离散赋值环的谱这一事实，\(G\setminus (G_y\setminus  G_y^0)\) 是 \(G\) 的一个开集，\(G\to  S\) 在其上诱导开态射，因而在 \(G_y^0\) 的每一点处给出“拟截面”。所以，必要时再扩张底 \(S\)，可以假设存在 \(G^0\) 在 \(S\) 上的截面 \(g'_2\)，使 \(g'_2(y)=g_2(y)\)。令 \(g'=g_1g'_2\)；由构造，\(g'(y)=g(y)=x\) 且 \(g'(y')=g_1(y')g'_2(y')\in  G^\sigma \)，所以 \(g(y')=x'\) 是 \(x\) 在 \(G^\sigma \) 中位于 \(y'\) 上方的泛化。这就证明了 (iii)。\par{}最后证明断言 (iv)。只须证明，若 \(x\in  G^\rho \)，则 \(x\) 的每个泛化 \(x'\) 都属于 \(G^\rho \)。（对适当的 \(h\)，取在 \(\varphi _{n^h}\) 下的像后，）可以假设甚至有 \(x\in  G^0\)。于是，对于每个与特征互素的整数 \(n>0\)，\(x\) 都在 \(\varphi _n\) 的像中（因为在特征与 \(n\) 互素的域上的连通有限型群中，\(n\) 次幂映射是满射）。由于 \(\varphi _n\) 为开映射，\(x'\) 也在 \(\varphi _n\) 的像中。更确切地说，取 \(G^\tau \) 中包含 \(G_y^0\) 的拟紧开集 \(U\)，则对所有与 \(p\) 互素的 \(n\)，都有 \(x'\in \varphi _n(U)\)。取 \(n\) 为各 \(\nu (z)\)（\(z\in  U\)）的与 \(p\) 互素的因子的公倍数，便得 \(x'\in  G^\rho \)。\par{}\hyperref[fga3.vi-1-theorem-1.1]{定理~\ref*{fga3.vi-1-theorem-1.1}} 的断言 (iii) 还有如下补充：\begin{corollary}{1.5}\setnumber{1.5}\label{fga3.vi-1-corollary-1.5}\skipsinglepar\par{}设 \(n>1\) 是一个整数，使 \(n\) 次幂同态 \(\varphi _n\colon  G\to  G\) 普遍开（例如 étale）。令 \(G^{(n)}=\varphi _n^{-1}(G^0)\)，并假设连通纤维 \(G_s^0\) “不含加法分量”（即将 \(k(s)\) 扩张到代数闭包所得的群不含同构于 \(\operatorname {G_a}\) 的子群）。那么，\(G\to  S\) 在 \(G^{(n^h)}\) 的各点处普遍开。特别地，若各 \(G_s^0\) 不含加法分量，且对每个整数 \(n>1\)，同态 \(\varphi _n\) 在剩余特征与 \(p\) 互素的各点处普遍开，则 \(G\to  S\) 在 \(G^\sigma \) 的各点处普遍开；从而（\hyperref[fga3.vi-1-theorem-1.1]{定理~\ref*{fga3.vi-1-theorem-1.1}} 的 (iii)），若 \(G^0\) 为闭集，则 \(G^\sigma \to  S\) 普遍开。在这种情形中，\(s\mapsto \dim  G_s\) 是 \(S\) 上的局部常值函数。\begin{proof}\par{}由假设，\(\varphi _n\) 的核 \({}_nG\) 在 \(S\) 上普遍开，因此 \(G\to  S\) 在 \({}_nG\) 的各点处普遍开，进而（将 \(n\) 换为 \(n^h\)）在 \({}_{(n^h)}G\) 的各点处也普遍开。但是，关于纤维 \(G_s^0\) 的假设恰好说明，\(G_s\) 中阶为 \(n\) 的某个幂的点在各 \(G^{(n^h)}\) 的并中稠密。容易由此得到，\(G\to  S\) 在此并集的所有点处普遍开，特别地，在单位截面沿线普遍开；又容易由此得到 \(s\mapsto \dim  G_s\) 为局部常值函数。\end{proof}\end{corollary}\begin{remark}{1.6}\setnumber{1.6}\label{fga3.vi-1-remark-1.6}\skipsinglepar\par{}回忆，态射 \(X\to  S\) 称为普遍开的，若它将每个开集映为开集，且任意基变换 \(S'\to  S\) 后仍保有这一性质。这也等价于下述条件（当 \(S\) 局部 Noether、\(X\to  S\) 局部有限型时），\(X\) 的每个不可约分支都支配 \(S\)，且此性质在任意基变换 \(S'\to  S\) 后保持。在这两个断言中，由于上述有限性假设，只须对这样的基变换 \(S'\to  S\) 验证：\(S'\) 是离散赋值环的谱（愿意的话，还可要求完备且剩余域代数闭……）。此定义显然推广到 \(X\) 的子集 \(Z\)（例如 \(G\) 的子集 \(G^\sigma \)）。即使从几何纤维连通的光滑射影态射 \(X\to  S\) 出发（例如 \(S=\operatorname {Spec}\mathbb {C}[j]\) 上椭圆曲线模族的纤维平方），\(\underline {\operatorname {Pic}}_{X/S}\) 在 \(S\) 上也可能并不普遍开，即 \(\underline {\operatorname {Pic}}_{X/S}\) 可能有完全位于 \(S\) 的某一个点上方的不可约分支；这是完全正常的现象。它与 \(X/S\) 的纤维的 Néron–Severi 群的秩可能向上跳跃这一事实有关（“复乘”现象）。不过，(1.5) 保证，在良好情形中，\(\underline {\operatorname {Pic}}_{X/S}^\sigma \)（而且通常看来甚至 \(\underline {\operatorname {Pic}}_{X/S}^\tau \)）在 \(S\) 上普遍开。\end{remark}\par{}最后给出一个有用的情形，在这一特殊情形下，\(G^0\) 的确是开集：\begin{corollary}{1.7}\setnumber{1.7}\label{fga3.vi-1-corollary-1.7}\skipsinglepar\par{}假设 \(G\to  S\) 在 \(G^0\) 的各点处普遍开（参见 \hyperref[fga3.vi-1-corollary-1.5]{推论~\ref*{fga3.vi-1-corollary-1.5}}），且纤维 \(G_s\) 可分，因而在 \(k(s)\) 上光滑（由 Cartier 的一个结果，在剩余特征为零时，后一个条件自动满足）。那么 \(G^0\) 在 \(G\) 中开。若 \(S\) 还约化，则 \(G^0\) 在 \(S\) 上光滑（特别地，平坦）。\begin{proof}\par{}第一个断言可进一步精确表述如下：对某个给定的 \(s\in  S\)，若 \(G\to  S\) 在 \(G_s^0\) 的各点处普遍开，且 \(G_s^0\) 在 \(k(s)\) 上可分，则 \(G^0\) 是 \(G_s^0\) 的一个邻域（而且对所有邻近点，关于 \(s\) 所作的假设仍然满足）。这一陈述的一个证明见 [\hyperref[GD1960]{2}，IV，§7]；此证明与群结构完全无关。\hyperref[fga3.vi-1-corollary-1.7]{推论~\ref*{fga3.vi-1-corollary-1.7}} 的最后一个断言也与群结构及连通分支问题无关；它是 [IV]，第5节 中一个平坦性判据的特例。更一般地，该判据给出：\end{proof}\end{corollary}\begin{corollary}{1.8}\setnumber{1.8}\label{fga3.vi-1-corollary-1.8}\skipsinglepar\par{}设 \(U\) 是纤维 \(G_s\) 的一个开子集，使得 \(G\to  S\) 在 \(U\) 的各点处普遍开（参见 (1.5)）。若 \(G_s\) 在 \(k(s)\) 上可分，且 \(S\) 在 \(s\) 处约化，则 \(G\) 在 \(U\) 的各点处在 \(S\) 上平坦（从而在 \(U\) 的各点处也在 \(S\) 上光滑）。\end{corollary}\begin{remarks}{1.9}\setnumber{1.9}\label{fga3.vi-1-remarks-1.9}\skipsinglepar\par{}若 \(G\) 在 \(S\) 上分离，我不知道 \(G^0\) 或 \(G^\tau \) 是否总在 \(G\) 中闭；这似乎不大可能。无论如何，去掉分离假设就有明显的反例，例如取具有可数无穷多个原点副本的仿射直线，便得到仿射直线上的群概形 \(G\)，其所有纤维都是平凡群，唯有一个纤维是 \(\mathbb {Z}\)。此群概形是一个 Picard 预概形的开群子概形，亦即单位截面的闭包；这里所说的 Picard 预概形对应于 \(S\)-概形 \(X\)，后者是一族退化为两条相交直线的圆锥曲线。\par{}即使从离散赋值环 \(V\) 的谱 \(S\) 上的有限平坦群概形 \(G\) 出发，\(G^0\) 化为单位截面，因而为闭集，只要去掉某些假设，各种结论仍会失效。假设 \(V\) 等特征且特征为 \(p>0\)，令 \(G\) 为函子同态 \(f\mapsto  f^p-tf\) 所给出的同态 \(\operatorname {G_a}\to \operatorname {G_a}\) 的核，其中 \(t\) 是 \(V\) 的一个素元。（回忆，按照定义，\(S\) 上的“加法群”\(\operatorname {G_a}\) 表示函子 \(S'\mapsto \Gamma (S',\mathcal {O}_{S'})\)。）于是，作为 \(S\)-概形，\(G=\operatorname {Spec} V[x]/(x^p-tx)\) 是 \(p\) 条相交“直线”的并，给出一个退化为无穷小加法型群的群 \(\mathbb {Z}/p\mathbb {Z}\)。在这个例子中，\(G^0=G^\sigma \)（这是 \(p\) 条直线中的一条），所以 \(G^0\) 在 \(G\) 中不开。在不等特征、剩余特征为 \(p>0\) 的情形中，从群概形 \(H=\mu _p\) 出发，它由 \(\operatorname {G_m}\) 中 \(p\) 次幂映射的核给出，故 \(H=\operatorname {Spec} V[x]/(x^p-1)\)；这同样是 \(p\) 条相交直线的并，给出群 \(\mathbb {Z}/p\mathbb {Z}\)（在特征 \(0\) 时），它退化为特征 \(p\) 的无穷小乘法型群。令 \(H'\) 为通常的有限群 \(\mathbb {Z}/p\mathbb {Z}\) 所定义的“常值群概形”，它是 \(\mathbb {Z}/p\mathbb {Z}\) 个 \(S\) 的副本的不交和，也就是 \(\operatorname {Spec} V^{\mathbb {Z}/p\mathbb {Z}}\)。于是，\(G=H\times _S H'\) 描述了特征 \(0\) 下的群 \((\mathbb {Z}/p\mathbb {Z})^2\)，它退化为特征 \(p\) 下的一个无穷小群与 \(\mathbb {Z}/p\mathbb {Z}\) 的乘积。这里，\(G^\rho \) 是单位截面与特殊纤维之并，故 \(G^\rho \) 不开；而 \(G^\sigma \) 是单位截面与泛纤维之并，故不闭。这与 \hyperref[fga3.vi-1-theorem-1.1]{定理~\ref*{fga3.vi-1-theorem-1.1}} 的 (iii)、(iv) 中等特征时成立的结论不同。当然，这些都是与特征 \(p>0\) 有关的现象。上述结果给出：\end{remarks}\begin{corollary}{1.10}\setnumber{1.10}\label{fga3.vi-1-corollary-1.10}\skipsinglepar\par{}假设 \(S\) 的剩余特征均为 \(0\)，于是 \(G^\tau =G^\sigma \) 且 \(G^0=G^\rho \)。那么 \(G^\tau =G^\sigma \) 为开集；若 \(G\) 在 \(S\) 上分离，且各 \(G_s^0\) 固有，则它甚至既开又闭。在相同假设下，\(G^0=G^\rho \) 在 \(S\) 上固有，因而在 \(G\) 中闭。最后，若对每个整数 \(n>1\)，\(G\) 中的 \(n\) 次幂同态普遍开，则 \(G^0\) 为开集。若各 \(G_s^0\) 还没有加法分量（例如，像上面一样各 \(G_s^0\) 固有），则 \(G^\tau \to  S\) 普遍开；若 \(S\) 约化，它甚至是光滑的。\end{corollary}\par{}最后指出下述容易得到的结果：\begin{proposition}{1.11}\setnumber{1.11}\label{fga3.vi-1-proposition-1.11}\skipsinglepar\par{}\(S\) 中存在开子集 \(U\)，使得 \(G\) 中满足 \(G\) 在 \(S\) 上光滑（分别地，平坦）的点所成的集合 \(G'\)，是由 \(G|U\) 诱导的某个开群子概形的底层集合。此外，\(G\) 在 \(U\) 上的每个截面都是 \(G'\) 在 \(U\) 上的截面。\end{proposition}\begin{corollary}{1.12}\setnumber{1.12}\label{fga3.vi-1-corollary-1.12}\skipsinglepar\par{}若 \(G\) 在单位截面的各点处在 \(S\) 上光滑（分别地，平坦），则它在 \(G\) 在 \(S\) 上的每个截面的各点处，以及在 \(G^0\) 的所有点处，也都是光滑的（分别地，平坦的）。若进一步对每个整数 \(n>0\)，\(n\) 次幂同态 \(\varphi _n\colon  G\to  G\) 在特征与 \(n\) 互素的各点处 étale，则 \(G\) 在 \(G^\sigma \) 的所有点处在 \(S\) 上光滑（分别地，平坦）。\end{corollary}\chapter{对 Picard 概形的局部性质的应用}\thispagestyle{fancy}\label{fga3.vi-2}\begin{theorem}{2.1}\setnumber{2.1}\label{fga3.vi-2-theorem-2.1}\skipsinglepar\begin{enumerate}
    
 \item[i.]设 \(f\colon  X\to  S\) 是固有光滑态射，并假设 \(\underline {\operatorname {Pic}}_{X/S}\) 存在（例如 \(f\) 射影时）。那么，\(\underline {\operatorname {Pic}}_{X/S}\) 在 \(S\) 上分离；而且对 \(\underline {\operatorname {Pic}}_{X/S}\) 的每个在 \(S\) 上有限型的闭子集 \(Z\)，都有 \(Z\) 在 \(S\) 上固有。

    
 \item[ii.]设 \(X\) 是域 \(k\) 上固有且几何正规的预概形。那么 \(\underline {\operatorname {Pic}}_{X/S}^0\) 在 \(k\) 上固有。

  \end{enumerate}\begin{proof}\par{}对于 (i)，由 [\hyperref[GD1960]{2}，II，§7] 的赋值判据，只须证明：若 \(S\) 是完备离散赋值环的谱，\(U\) 是由 \(S\) 的泛点组成的开集，则 \(\underline {\operatorname {Pic}}_{X/S}\) 在 \(S\) 上的每个有理截面，即在 \(U\) 上的每个截面，都唯一延拓为 \(S\) 上的截面。考虑到 \(\underline {\operatorname {Pic}}_{X/S}\) 的定义，这等价于下述陈述：对 \(V=f^{-1}(U)\) 上的每个可逆模 \(\mathscr {L}\)，\(X\) 上存在一个延拓 \(\mathscr {L}\) 的可逆模，且在同构意义下唯一。用“Cartier”除子类来描述 \(V\) 上（分别地，\(X\) 上）的可逆模，很容易得到这一点：\(X\) 的局部环是正则的（因为 \(X\) 在正则的 \(S\) 上光滑），因而由 Auslander 的结果，它们都是析因环，所以 \(S\) 上的每个除子都是 Cartier 除子。事实上，\(V\) 上的每个除子，取其“闭包”即可延拓为 \(X\) 上的除子。\par{}对于 (ii)，利用 Chow 引理可归约到 \(X\) 射影、因而嵌入某个 \(\mathbb {P}_k^n\) 的情形；还可假设 \(X\) 连通。若 \(\dim  X=1\)，则 \(X\) 在 \(k\) 上光滑，可以应用 (i)。若 \(\dim  X\geqslant 2\)，则可以利用已知“等价判据”的一个变式：存在有限条在 \(X\) 上光滑的曲线 \(Y_i\)（由 \(X\) 与 \(\mathbb {P}_k^n\) 中适当线性子空间相交而得），使得 \(\underline {\operatorname {Pic}}_{X/k}^\tau \to \prod _i\underline {\operatorname {Pic}}_{Y_i/k}^\tau \) 为单态射，从而在连通分支上也诱导单态射。由上述讨论，靶在 \(k\) 上固有；这里的态射是群概形同态，必为闭浸入，因此 \(\underline {\operatorname {Pic}}_{X/k}^0\) 在 \(k\) 上也固有。利用代数闭域上交换代数群的结构（Chevalley–Borel 的结果），可以避免使用微妙的等价判据。此时只须证明，从去掉原点的仿射直线到 \(\underline {\operatorname {Pic}}_{X/S}^\tau \) 的每个态射都是常值态射；这等价于说，\(X[t,t^{-1}]\) 上的每个可逆模都来自 \(X\) 上的可逆模。这是一个熟知的初等结果，甚至不需要 \(X\) 在 \(k\) 上固有这一事实（\(X\) 正规的假设使我们可立即归约到 \(X\) 正则的情形）。\end{proof}\end{theorem}\begin{remark}{2.2}\setnumber{2.2}\label{fga3.vi-2-remark-2.2}\skipsinglepar\par{}即使仅假设 \(f\) 平坦，且其纤维 \(X_s\) 为局部完全交并在余维 \(\leqslant 2\) 处在 \(k(s)\) 上光滑，上面对 (i) 的证明仍然成立；这里利用了 \hyperref[fga3.vi-bib-5]{[5]} 中证明的下述事实：在余维 \(\leqslant 3\) 处正则的 Noether 完全交局部环是析因环（“Samuel 猜想”）。注意，若将“余维 \(\leqslant 2\)”换为“余维 \(\leqslant 1\)”，即换为“正规”的假设，结论就不再成立；考虑一族退化为二次锥面的非奇异二次曲面即可看出。\end{remark}\begin{corollary}{2.3}\setnumber{2.3}\label{fga3.vi-2-corollary-2.3}\skipsinglepar\par{}设 \(f\colon  X\to  S\) 是固有正规态射（即平坦且几何纤维正规），并假设 \(\underline {\operatorname {Pic}}_{X/S}\) 存在。那么 \(\underline {\operatorname {Pic}}_{X/S}^0\) 在 \(S\) 上固有，从而在 \(\underline {\operatorname {Pic}}_{X/S}\) 中闭；此外，\(\underline {\operatorname {Pic}}_{X/S}^\tau \) 与 \(\underline {\operatorname {Pic}}_{X/S}^\rho \) 都闭，在“等特征”情形下，\(\underline {\operatorname {Pic}}_{X/S}^\sigma \) 也闭。\begin{proof}\par{}应用 \hyperref[fga3.vi-1-theorem-1.1]{定理~\ref*{fga3.vi-1-theorem-1.1}} 及 \hyperref[fga3.vi-2-theorem-2.1]{定理~\ref*{fga3.vi-2-theorem-2.1}} 的 (ii)。\end{proof}\end{corollary}\begin{corollary}{2.4}\setnumber{2.4}\label{fga3.vi-2-corollary-2.4}\skipsinglepar\par{}设 \(f\colon  X\to  S\) 是固有光滑态射，且 \(\underline {\operatorname {Pic}}_{X/S}\) 存在，并为 \(S\) 上有限型概形 \(P^{(i)}\) 的和（参见 \hyperref[fga3.v-4-proposition-4.1]{[3]，V，4.1}）。那么每个 \(P^{(i)}\) 都在 \(S\) 上固有。\begin{proof}\par{}这由 \hyperref[fga3.vi-2-theorem-2.1]{定理~\ref*{fga3.vi-2-theorem-2.1}} 的 (i) 推出。\end{proof}\end{corollary}\par{}如 \hyperref[fga3.vi-2-remark-2.2]{注~\ref*{fga3.vi-2-remark-2.2}} 中所指出，减弱对 \(f\) 的纤维所作的假设，仍可推广上述结果；但若仅假设 \(f\) 正规，则结果不再成立。在这种情形中，即使假定 \(\underline {\operatorname {Pic}}_{X/S}^\tau \) 在 \(S\) 上有限型，我也不知道它是否仍在 \(S\) 上固有。\begin{theorem}{2.5}\setnumber{2.5}\label{fga3.vi-2-theorem-2.5}\skipsinglepar\par{}设 \(f\colon  X\to  S\) 是固有平坦态射，且 \(\underline {\operatorname {Pic}}_{X/S}\) 存在；对每个整数 \(n\)，以 \(\varphi _n\) 表示此群预概形中的 \(n\) 次幂同态。那么，\(\varphi _n\) 在剩余特征与 \(n\) 互素的所有点 \(x\in  X\) 处 étale。\end{theorem}\par{}由 étale 态射的无穷小刻画，上面的断言等价于下述陈述：\begin{lemma}{2.6}\setnumber{2.6}\label{fga3.vi-2-lemma-2.6}\skipsinglepar\par{}假设 \(S\) 是 Artin 局部环 \(A\) 的谱，其极大理想 \(\mathfrak {m}\) 的 \((\nu +1)\) 次幂为零，并令 \(A_{\nu -1}=A/\mathfrak {m}^\nu \)、\(X_{\nu -1}=X\otimes _A A_{\nu -1}\)。设 \(\mathscr {L}\) 是 \(X\) 上的可逆模，\(\mathscr {L}'_{\nu -1}\) 是 \(X_{\nu -1}\) 上的可逆模，其 \(n\) 次张量幂同构于 \(\mathscr {L}_{\nu -1}=\mathscr {L}\otimes _A A_{\nu -1}\)。那么，\(X\) 上存在可逆模 \(\mathscr {L}'\)，其 \(n\) 次张量幂同构于 \(\mathscr {L}\)（只要 \(n\) 与 \(k=k(\mathscr {L})\) 的剩余特征互素）。\begin{proof}\par{}令 \(V=\mathfrak {m}^\nu =\mathfrak {m}^{\nu }/\mathfrak {m}^{\nu +1}\)，它是 \(k=k(A)\) 上的向量空间。先将 \(\mathscr {L}'_{\nu -1}\) 任意延拓为 \(X\) 上的可逆模 \(\mathscr {L}'\)。这样做的障碍位于 \(\operatorname {H}^2(X_0,\mathcal {O}_{X_0})\otimes _k V\) 中；但由关于 \(\mathscr {L}'_{\nu -1}\) 的假设及 \(\mathscr {L}_{\nu -1}\) 可以延拓这一事实，此障碍与 \(n\) 的乘积为零；由于 \(n\) 与特征互素，障碍本身必为零。延拓的不唯一性位于 \(\operatorname {H}^1(X_0,\mathcal {O}_{X_0})\otimes _k V\) 中，而 \({\mathscr {L}'}^{\otimes  n}\) 与 \(\mathscr {L}\) 之间的偏差 \(\xi \) 也位于同一个模中。若想修正 \(\mathscr {L}'\) 使此偏差为零，就须在上述模中找到一个 \(\eta \)，使 \(n\eta =\xi \)。由于 \(n\) 与特征互素，这仍然可以做到。\end{proof}\end{lemma}\begin{corollary}{}\label{fga3.vi-2-corollary}\skipsinglepar\par{}在 \hyperref[fga3.vi-2-theorem-2.5]{定理~\ref*{fga3.vi-2-theorem-2.5}} 的条件下，进一步假设纤维 \(X_s\) 的 Picard 概形均不含加法分量（例如，各 \(X_s\) 几何正规时；参见 \hyperref[fga3.vi-2-theorem-2.1]{定理~\ref*{fga3.vi-2-theorem-2.1}} 的 (ii)）。那么 \(\underline {\operatorname {Pic}}_{X/S}\to  S\) 在 \(\underline {\operatorname {Pic}}_{X/S}^\sigma \) 的各点处普遍开。若 \(\underline {\operatorname {Pic}}_{X/S}^0\) 闭（例如，各 \(X_s\) 几何正规时；参见 \hyperref[fga3.vi-2-corollary-2.3]{推论~\ref*{fga3.vi-2-corollary-2.3}}），则 \(\underline {\operatorname {Pic}}_{X/S}^\sigma \) 本身在 \(S\) 上普遍开。最后，在等特征情形中，\(\underline {\operatorname {Pic}}_{X/S}^\rho \to  S\) 普遍开。\begin{proof}\par{}应用 \hyperref[fga3.vi-1-corollary-1.5]{推论~\ref*{fga3.vi-1-corollary-1.5}} 和 \hyperref[fga3.vi-1-theorem-1.1]{定理~\ref*{fga3.vi-1-theorem-1.1}}。\end{proof}\end{corollary}\begin{corollary}{2.7}\setnumber{2.7}\label{fga3.vi-2-corollary-2.7}\skipsinglepar\par{}设 \(f\colon  X\to  Y\) 是固有平坦态射，且 \(\underline {\operatorname {Pic}}_{X/S}\) 存在。那么，\(S\) 上的函数 \(s\mapsto \dim \underline {\operatorname {Pic}}_{X_s/k(s)}\) 上半连续（即只可能向上跳跃，不能向下跳跃）；若各 \(\underline {\operatorname {Pic}}_{X_s/k(s)}\) 不含加法分量，则它甚至连续（即局部常值）。\begin{proof}\par{}对于局部 Noether 底上任意局部有限型群预概形，第一个断言都是平凡的，或几乎平凡的，因为只须考察单位截面沿线的情形。第二个断言由 \hyperref[fga3.vi-2-theorem-2.5]{定理~\ref*{fga3.vi-2-theorem-2.5}} 得到。\end{proof}\end{corollary}\begin{remark}{2.8}\setnumber{2.8}\label{fga3.vi-2-remark-2.8}\skipsinglepar\par{}设 \(s,s'\in  S\)，且 \(s\) 是 \(s'\) 的特殊化。那么，\hyperref[fga3.vi-2-corollary-2.7]{推论~\ref*{fga3.vi-2-corollary-2.7}} 等价于 \(X_{s'}\) 及其“特殊化”\(X_s\) 的 Picard 概形的维数之间的一个不等式（分别地，等式）。甚至在 Picard 概形构造出来以前，Serre 就已注意到：对于 \emph{光滑} 态射 \(f\colon  X\to  S\)，各 \(X_s\) 的 Picard 簇的维数不变这一事实，是基本群特殊化理论（[\hyperref[Gro1960b]{4}，X]）以及 Picard 簇的有限阶点与基本群 Abel 化之间的经典 Kummer 型关系（[\hyperref[Gro1960b]{4}，XI]）的形式推论。用 \(\alpha \)、\(\mu \)、\(\lambda \) 分别表示 \(\underline {\operatorname {Pic}}_{X_s/k(s)}\) 的 Abel 部分、乘法部分和加法部分的维数，并类似定义 \(\alpha '\)、\(\mu '\)、\(\lambda '\)，则已知关系可写为下列不等式：\begin{eqn}{*}\setnumber{*}\label{fga3.vi-2-equation-star}\skipsinglepar\[         \alpha +\mu +\lambda          \geqslant  \alpha '+\mu '+\lambda '       \tag{*}       \]\end{eqn}（只要 \(\underline {\operatorname {Pic}}_{X/S}\) 存在就成立，因此很可能在所有情形中均成立）。当 \(\lambda =\lambda '=0\) 时，这个不等式化为在同样存在性假设下成立的等式：\[       \alpha +\mu        = \alpha '+\mu '.     \]。由 Serre 的论证，若各 \(X_s\) 可分（甚至不必假设 \(\underline {\operatorname {Pic}}_{X/S}^\tau \) 存在），或者各 \({}_n\underline {\operatorname {Pic}}_{X/S}\)（Picard 函子中各 \(\varphi _n\) 的核）在 \(S\) 上 \emph{分离}，则还有 \begin{eqn}{**}\setnumber{**}\label{fga3.vi-2-equation-star-star}\skipsinglepar\[         2\alpha +\mu          \leqslant  2\alpha '+\mu '       \tag{**}       \]\end{eqn}（这里利用了 \hyperref[fga3.vi-2-theorem-2.5]{定理~\ref*{fga3.vi-2-theorem-2.5}} 所保证的它们在 \(S\) 上 étale 这一事实）。我们倾向于猜想，\hyperref[fga3.vi-2-equation-star]{式~\ref*{fga3.vi-2-equation-star}} 在所有情形中均为等式，至少在各 \(X_s\) 可分时如此；还猜想有不等式 \begin{eqn}{***}\setnumber{***}\label{fga3.vi-2-equation-star-star-star}\skipsinglepar\[         \begin {aligned}           \alpha  &\leqslant  \alpha '         \\\lambda  &\geqslant  \lambda '         \end {aligned}       \tag{***}       \]\end{eqn}。只要所考虑的是局部 Noether 的 \(S\) 上的局部有限型群预概形，且纤维维数恒定，这些不等式就应当成立（这个方向的一个肯定结果见 \hyperref[fga3.vi-1-lemma-1.3]{引理~\ref*{fga3.vi-1-lemma-1.3}}）。\end{remark}\begin{remark}{2.9}\setnumber{2.9}\label{fga3.vi-2-remark-2.9}\skipsinglepar\par{}在所有已知情形中，\(\underline {\operatorname {Pic}}_{X/S}^\tau \) 在 \(S\) 上普遍开；但即使 \(f\colon  X\to  S\) 光滑，我们恐怕也不宜过于乐观。无论如何，Mumford 通过对 Igusa 曲面作无穷小变形，构造了 \(\underline {\operatorname {Pic}}_{X/S}^\tau \) 在 \(S\) 上不平坦的例子（诚然，其中 \(S\) 不约化，实际上 \(S\) 是 Artin 环的谱）。Mumford 所考察的点也属于 \(\underline {\operatorname {Pic}}_{X/S}^\rho \)；因此，（对光滑的 \(f\colon  X\to  S\)）\(\underline {\operatorname {Pic}}_{X/S}\) 在 \(\underline {\operatorname {Pic}}_{X/S}^\sigma \) 的各点处在 \(S\) 上平坦仍有可能成立。不过，报告人怀疑这一点并不总成立，即使只考虑 \(\underline {\operatorname {Pic}}_{X/S}^0\) 的各点，并假设 \(S\) 是离散赋值环的谱。这个问题与研究 Abel 概形在有限自同构群作用下的不动点有关；在这种情形中，我们似乎还缺少例子。看来，即使只考虑光滑射影的 \(f\colon  X\to  S\)，如果不对 \(f\) 的纤维的性质作更特殊的假设，本节关于 \(\underline {\operatorname {Pic}}_{X/S}\) 的局部正则性的结果以及 \hyperref[fga3.vi-2-remark-2.8]{注~\ref*{fga3.vi-2-remark-2.8}} 中提出的猜想，基本上已经穷尽了在这方面所能说的内容。不过，回忆，若 \(\underline {\operatorname {Pic}}_{X/S}\) 的几何纤维 \emph{约化} 且没有加法分量，则由 \hyperref[fga3.vi-1-corollary-1.8]{推论~\ref*{fga3.vi-1-corollary-1.8}} 和 \hyperref[fga3.vi-2-theorem-2.5]{定理~\ref*{fga3.vi-2-theorem-2.5}} 可知，只要 \(S\) 约化，\(\underline {\operatorname {Pic}}_{X/S}\) 就在 \(\underline {\operatorname {Pic}}_{X/S}^\sigma \) 的各点处在 \(S\) 上光滑。若 \(f\colon  X\to  S\) 正规，则即使不对 \(S\) 作任何假设，此结果仍成立；这将在 \hyperref[fga3.vi-3-theorem-3.5]{定理~\ref*{fga3.vi-3-theorem-3.5}} 中看到。就此指出：\end{remark}\begin{proposition}{2.10}\setnumber{2.10}\label{fga3.vi-2-proposition-2.10}\skipsinglepar\begin{enumerate}
    
 \item[i.]若 \(\underline {\operatorname {Pic}}_{X/S}\) 在单位截面的各点处在 \(S\) 上光滑（分别地，平坦），则它在 \(\underline {\operatorname {Pic}}_{X/S}^\sigma \) 的所有点处，以及 \(\underline {\operatorname {Pic}}_{X/S}\) 在 \(S\) 上的每个截面的各点处，也都是光滑的（分别地，平坦的）。

    
 \item[ii.]设 \(s\in  S\) 满足 \(\operatorname {H}^2(X_s,\mathcal {O}_{X_s})=0\)。那么 \(s\) 有一个开邻域 \(U\)，使 \(\underline {\operatorname {Pic}}_{X/S}|U\) 在 \(U\) 上光滑。

    
 \item[iii.]设 \(X\) 是一个域上的固有概形。那么有 \[         \dim \underline {\operatorname {Pic}}_{X/k}         \leqslant  \dim \operatorname {H}^1(X,\mathcal {O}_X)       \]，其中等号成立当且仅当 \(\underline {\operatorname {Pic}}_{X/k}\) 在 \(k\) 上光滑；若 \(k\) 的特征为 \(0\)，则总是如此。

  \end{enumerate}\begin{proof}\par{}断言 (i) 由 \hyperref[fga3.vi-2-theorem-2.5]{定理~\ref*{fga3.vi-2-theorem-2.5}} 和 \hyperref[fga3.vi-1-corollary-1.12]{推论~\ref*{fga3.vi-1-corollary-1.12}} 得到；(ii) 则由光滑态射的无穷小判据及一个熟知的障碍计算得到，这里利用了关于 \(s\) 的假设由“半连续性定理”对邻近点仍成立。最后，(iii) 由 \(\operatorname {H}^1(X,\mathcal {O}_X)\) 同构于 \(\underline {\operatorname {Pic}}_{X/k}\) 在单位元处的 Zariski 切空间这一事实得到；最后一个断言是 Cartier 定理的特例，该定理说，特征 \(0\) 下的一个“形式群”在 \(k\) 上形式光滑。\end{proof}\end{proposition}\chapter{\texorpdfstring{\(\underline {\operatorname {Pic}}_{X/S}\)}{Pic(X/S)} 的典范 Abel 子概形与 Albanese 概形}\thispagestyle{fancy}\label{fga3.vi-3}\begin{proposition}{3.1}\setnumber{3.1}\label{fga3.vi-3-proposition-3.1}\skipsinglepar\par{}设 \(k\) 是域，\(G\) 是 \(k\) 上交换且“不含加法分量”的有限型群概形。那么，\(G_\mathrm {red}^0\) 在 \(k\) 上可分，因而是 \(k\) 上的光滑群概形。\begin{proof}\par{}若 \(k\) 是完美域，断言显然成立；特别地，对 \(G_{\overline {k}}\) 也成立，其中 \(\overline {k}\) 是 \(k\) 的代数闭包。因此，只须证明 \((G_{\overline {k}}^0)_\mathrm {red}\) 来自 \(G\) 的一个子概形。由 \(G_{\overline {k}}\) 不含加法分量的假设，容易得到存在整数 \(m\)，使 \((G_{\overline {k}}^0)_\mathrm {red}\) 是 \(G_{\overline {k}}\) 中 \(m\) 次幂同态的“概形论”像。由于后一个同态来自 \(G\) 的类似同态，该同态的概形论像就是所需对象。\end{proof}\end{proposition}\begin{corollary}{3.2}\setnumber{3.2}\label{fga3.vi-3-corollary-3.2}\skipsinglepar\par{}设 \(X\) 是 \(k\) 上的正规固有概形，且 \(\underline {\operatorname {Pic}}_{X/k}\) 存在。那么 \(\underline {\operatorname {Pic}}_{X/k}\) 有一个底层集合为 \(\underline {\operatorname {Pic}}_{X/k}^0\) 的 Abel 子概形。\begin{proof}\par{}由 \hyperref[fga3.vi-2-theorem-2.1]{定理~\ref*{fga3.vi-2-theorem-2.1}} 的 (ii)，\(\underline {\operatorname {Pic}}_{X/k}^0\) 在 \(k\) 上固有，因而满足 \hyperref[fga3.vi-3-proposition-3.1]{命题~\ref*{fga3.vi-3-proposition-3.1}} 的条件。\end{proof}\end{corollary}\par{}上述结果表明，在某些情形中，经典的“Picard 簇”（在现今理论中即 \((\underline {\operatorname {Pic}}_{X/\overline {k}})_\mathrm {red}^0\)）“定义在 \(k\) 上”，无须假设域 \(k\) 完美。\par{}现在设 \(f\colon  X\to  S\) 是固有平坦的相对概形，为简单起见设 \(\mathcal {O}_S\xrightarrow {\sim } f_*(\mathcal {O}_X)\)，并假设 \(\underline {\operatorname {Pic}}_{X/S}\) 存在，且 \(\underline {\operatorname {Pic}}_{X/S}^0\) 在 \(S\) 上固有。对于 \hyperref[fga3.vi-3-theorem-3.3]{定理~\ref*{fga3.vi-3-theorem-3.3}} 的 (ii)，进一步假设 \(\underline {\operatorname {Pic}}_{X/S}\) 中存在一个包含 \(\underline {\operatorname {Pic}}_{X/S}^0\)、在 \(S\) 上拟射影的开集；如上所见，若 \(f\) 射影且几何纤维可分且不可约，此条件就得到满足。回忆，\(S\) 上的 \emph{Abel 概形} 是 \(S\) 上的群概形，它在 \(S\) 上固有且光滑，并且几何纤维连通。我们拟考察是否存在 \(\underline {\operatorname {Pic}}_{X/S}\) 的群子概形 \(A\)，使其为 Abel 概形且底层集合是 \(\underline {\operatorname {Pic}}_{X/S}^0\)。刚刚已看到，若 \(S\) 是域的谱，这样的 \(A\) 总是存在。在这里所考虑的一般情形中，目前可以说的是：\begin{theorem}{3.3}\setnumber{3.3}\label{fga3.vi-3-theorem-3.3}\skipsinglepar\par{}在上述条件下：\begin{enumerate}
      
 \item[i.]若 \(\underline {\operatorname {Pic}}_{X/S}\) 有一个底层集合为 \(\underline {\operatorname {Pic}}_{X/S}^0\) 的 Abel 子概形，则它唯一。因此，它的构造与基变换相容。

      
 \item[ii.]这样的 Abel 子概形存在的充要条件是：作每个 \(S'\) 为局部 Artin 概形的基变换 \(S'\to  S\) 后，它都存在。若 \(S\) 是局部环的谱，则甚至只须对满足 \(A_n=A/\mathfrak {m}^{n+1}\) 的各 \(S'=\operatorname {Spec}(A_n)\) 验证。若 \(S\) 约化，则也只须对各 \(S'\) 是离散赋值环的谱的情形验证（愿意的话，可取完备且剩余域代数闭）。

      
 \item[iii.]假设 \(A\) 存在，设 \(B=\operatorname {Alb}^0(X/S)\) 是其对偶 Abel 概形（即 \(B=\operatorname {Pic}_{A/S}^0\)；参见 \hyperref[fga3.vi-bib-7]{[7]}）。那么，可以典范地构造 \(B\) 的一个主齐性空间 \(P=\operatorname {Alb}^1(X/S)\)，以及一个 \(S\)-态射 \(X\to  P\)；它对于从 \(X\) 到准 Abel 概形（即 Abel 概形的主齐性空间）的 \(S\)-态射是泛的。\(\operatorname {Alb}^0(X/S)\)、\(\operatorname {Alb}^1(X/S)\)、\(X\to \operatorname {Alb}^1(X/S)\) 的构造均与基变换交换。

    \end{enumerate}\begin{proof}\par{}证明梗概如下：\begin{enumerate}\item[i.]这是交换群概形的 Abel 子概形所具有的一般刚性性质：若两个这样的子概形在点 \(s\in  S\) 处在集合意义下相同，则它们在 \(s\) 的整个连通分支上相同（\hyperref[fga3.vi-bib-7]{[7]}）。（这一结果推广了 Chow 的一个经典定理。）

      \item[ii.]利用 Hilbert 概形可见，将 \(S\) 上的每个 \(S'\) 映到 \((\underline {\operatorname {Pic}}_{X/S})\times _S S'\) 的典范 Abel 子概形所成的集合（其中至多一个元素）的函子，可由 \(S\) 上的有限型概形 \(T\) 表示。由 (i)，\(T\to  S\) 为单态射；由 \hyperref[fga3.vi-3-corollary-3.2]{推论~\ref*{fga3.vi-3-corollary-3.2}}，它是满射。说 \(\underline {\operatorname {Pic}}_{X/S}\) 有典范 Abel 子概形，等价于说 \(T\) 在 \(S\) 上有截面，也等价于说 \(T\to  S\) 为同构。这又等价于说 \(T\to  S\) 为 étale，或者，在 \(S\) 约化时，等价于说 \(T\to  S\) 固有。于是立刻得到 (ii)。

      \item[iii.]只用 \(\underline {\operatorname {Pic}}_{X/S}\) 的定义，就可注意到：对于 \(S\) 上的每个 Abel 概形 \(C\)，给定从 \(X\) 到 \(C\) 的一个主齐性空间的 \(S\)-态射，等价于给定群同态 \(C'\to \underline {\operatorname {Pic}}_{X/S}\)，其中 \(C'\) 是 \(C\) 的对偶 Abel 概形。但若 \(\underline {\operatorname {Pic}}_{X/S}\) 的典范 Abel 子概形 \(A\) 存在，则这些同态必通过 \(A\) 分解（再次利用有限阶点即可看出）。于是立刻得到 (iii)。\qedhere\end{enumerate}\end{proof}\end{theorem}\begin{remark}{3.4}\setnumber{3.4}\label{fga3.vi-3-remarks-3.4}\skipsinglepar\par{}若 \(\underline {\operatorname {Pic}}_{X/S}\) 的典范 Abel 子概形存在，将其记为 \(\underline {\operatorname {Pic}}_{X/S}^{00}\)。遗憾的是，它并不总是存在；对 Igusa 曲面作无穷小变形（一阶模变形）即可看到这一点。不过，至少在 \(S\) 约化时，\(\underline {\operatorname {Pic}}_{X/S}^{00}\) 可能总存在；由 (ii)，这等价于在 \(S\) 为离散赋值环的谱时存在。设 \(X_0\) 和 \(X_1\) 分别为 \(X\) 的特殊纤维和泛纤维，令 \(A_1=\underline {\operatorname {Pic}}_{X_1/K_1}\)，其中 \(K\) 是赋值环 \(V\) 的分式域。由 Koizumi 的结果，\(S\) 上存在本质唯一的 Abel 概形 \(A\)，其泛纤维为 \(A_1\)。如同 (2.1), (i) 的证明那样（从现在起假设 \(X\) 在 \(S\) 上光滑），容易看出，\(A_1\) 的恒等态射延拓为态射 \[       A \to  \underline {\operatorname {Pic}}_{X/S}.     \]由此得到同态 \begin{eqn}{*}\setnumber{*}\label{fga3.vi-3-equation-star}\skipsinglepar\[         A_0 \to  \underline {\operatorname {Pic}}_{X_0/k}^{00}       .\tag{*}       \]\end{eqn}容易证明，此同态为满射，核是一个有限 \(p\)-群，其中 \(p\) 是剩余域 \(k\) 的特征指数（仍利用有限阶点）。于是，关于 \(X/S\) 的下列条件等价：\begin{enumerate}
      
 \item[a.]上述同态 \hyperref[fga3.vi-3-equation-star]{式~\ref*{fga3.vi-3-equation-star}} 是同构（Shimura 将此表述为“Picard 簇”的构造“与特殊化相容”）。

      
 \item[b.]\(\underline {\operatorname {Pic}}_{X/S}^{00}\) 存在（此时它恰为 \(A\)）。

      
 \item[c.]（为便于回忆。）各 \(\underline {\operatorname {Pic}}_{X_n/S}^{00}\) 存在。

    \end{enumerate}\par{}由上面对 \hyperref[fga3.vi-3-equation-star]{式~\ref*{fga3.vi-3-equation-star}} 的核所作的注记，剩余特征为零时，条件 (a) 得到满足；不过，\hyperref[fga3.vi-3-theorem-3.5]{定理~\ref*{fga3.vi-3-theorem-3.5}} 将显著推广这一结果。\par{}当然，若 \(\underline {\operatorname {Pic}}_{X/S}\) 在 \(\underline {\operatorname {Pic}}_{X/S}^0\) 的各点处在 \(S\) 上光滑，则后者在 \(\underline {\operatorname {Pic}}_{X/S}\) 中开（参见 \hyperref[fga3.vi-1-corollary-1.7]{推论~\ref*{fga3.vi-1-corollary-1.7}}）；赋予诱导结构后，它是 \(\underline {\operatorname {Pic}}_{X/S}\) 的一个 Abel 子概形，因而等于 \(\underline {\operatorname {Pic}}_{X/S}^{00}\)；后者在此情形中存在。但我们有更强的结果：\end{remark}\begin{theorem}{3.5}\setnumber{3.5}\label{fga3.vi-3-theorem-3.5}\skipsinglepar\par{}在 \hyperref[fga3.vi-3-theorem-3.3]{定理~\ref*{fga3.vi-3-theorem-3.3}} 的条件下，设 \(s\in  S\) 使得 \(\underline {\operatorname {Pic}}_{X_s/k(s)}\) 在 \(k(s)\) 上光滑（或者等价地，使得 \(\dim \underline {\operatorname {Pic}}_{X_s/k(s)}=\dim \operatorname {H}^1(X_s,\mathcal {O}_{X_s})\)）。那么，\(s\) 有一个开邻域 \(U\)，使 \(\underline {\operatorname {Pic}}_{X/S}\) 在 \(\underline {\operatorname {Pic}}_{X/S}^0|U\) 的各点处在 \(S\) 上光滑，因而后者是 \(\underline {\operatorname {Pic}}_{X/S}|U\) 中的开 Abel 子概形。特别地，\(\underline {\operatorname {Pic}}_{X|U/U}^{00}\) 存在。\begin{proof}\par{}说明证明的思路。上述结果使我们可归约到 \(S\) 是 Artin 局部环 \(A\) 的谱的情形，对 \(A\) 的无穷小阶数作归纳。因此可以假设 \(\underline {\operatorname {Pic}}_{X_n/A_n}^0\) 在 \(A_n\) 上光滑，只须证明 \(\underline {\operatorname {Pic}}_{X_{n+1}/A_{n+1}}^0\) 在 \(A_{n+1}\) 上光滑。为此，注意只须构造 \(A_{n+1}\) 上的 Abel 概形 \(P_{n+1}\) 来延拓 \(P_n=\underline {\operatorname {Pic}}_{X_n/A_n}^0\)，并构造 \(X_{n+1}\times _{A_{n+1}}P_{n+1}\) 上的可逆模 \(\mathscr {L}_{n+1}\)，来延拓 \(X_n\times _{A_n}P_n\) 上的可逆模 \(\mathscr {L}_n\)；后一个模出现于将 Picard 概形 \(\underline {\operatorname {Pic}}_{X_n/A_n}\) 定义为泛问题的解时。（注意：可以假设 \(X\) 在 \(S\) 上配有截面……。）为进行这一构造，必须使用下列关键结果：\emph{定义在 Artin 局部环的商环上的每个 Abel 概形都可以延拓}（换言之，代数闭域上的 Abel 概形的绝对“形式模概形”（\hyperref[fga3.ii]{[3]，II}）在 \(k\) 上的 Witt 向量环上光滑）。这一结果可由提升障碍的一般形式性质及群运算得到。利用此结果，先将 \(P_n\) 任意延拓为 \(P_{n+1}\)；然后得到提升 \(\mathscr {L}_n\) 的障碍，位于 \(\operatorname {H}^2(X_0\times  P_0,\mathcal {O}_{X_0\times  P_0})\otimes _k V\) 中（其中 \(V=\mathfrak {m}^{n+1}/\mathfrak {m}^{n+2}\)），更确切地说，位于子空间 \(\operatorname {H}^1(X_0,\mathcal {O}_{X_0})\otimes \operatorname {H}^1(P_0,\mathcal {O}_{P_0})\otimes _k V\) 中（因为 \(\mathscr {L}_n\) 在两个因子 \(X_n\) 和 \(P_n\) 上的限制平凡）。但后一个空间恰为 \(\operatorname {H}^1(P_0,\mathscr {G}_{P_0/k})\otimes  V\)，其中 \(\mathscr {G}_{P_0/k}\) 是 \(P_0/k\) 的切层；因而它也是将 \(P_n\) 提升为 \(P_{n+1}\) 时描述不确定性的空间（\hyperref[fga3.ii]{[3]，II}）。所以，可以修正此提升（而且正如预期的那样，修正方式唯一），使提升 \(\mathscr {L}_n\) 的障碍消失。\end{proof}\end{theorem}\begin{corollary}{3.6}\setnumber{3.6}\label{fga3.vi-3-corollary-3.6}\skipsinglepar\par{}在 \hyperref[fga3.vi-3-theorem-3.5]{定理~\ref*{fga3.vi-3-theorem-3.5}} 的条件下，\(\operatorname {R}^1 f_*(\mathcal {O}_X)\) 是 \(S\) 上在 \(s\) 的某个邻域内局部自由的模，其构造与基变换交换。\begin{proof}\par{}此模恰好是 \(\underline {\operatorname {Pic}}_{X/S}\) 沿单位截面的切模。\end{proof}\end{corollary}\begin{remark}{3.7}\setnumber{3.7}\label{fga3.vi-3-remark-3.7}\skipsinglepar\par{}利用与 \hyperref[fga3.vi-3-theorem-3.5]{定理~\ref*{fga3.vi-3-theorem-3.5}} 相同的论证，可以证明：若 \(S'\) 是 \(S\) 中由幂零凝聚理想定义的子概形，且仅假设 \(\underline {\operatorname {Pic}}_{X'/S'}^0\) 存在并具有光滑纤维，则 \(\underline {\operatorname {Pic}}_{X/S}^0\) 必然存在，且为 \(S\) 上的 Abel 概形。这使我们在不作任何射影性假设时，也能在某些情形中构造 Picard 概形；例如，Artin 环上 Abel 概形的对偶 Abel 概形总存在。在 \(X\) 是 \(S\) 上 Abel 概形的情形中，利用 \hyperref[fga3.vi-3-corollary-3.6]{推论~\ref*{fga3.vi-3-corollary-3.6}}，以及 \(\operatorname {H}^\bullet (X_s,\mathcal {O}_{X_s})\) 是 \(\operatorname {H}^1(X_s,\mathcal {O}_{X_s})\) 上外代数的已知结构（Rosenlicht–Serre），便得 \(\operatorname {R}^if_*(\mathcal {O}_X)\) 对 \emph{所有} \(i\) 均局部自由；更确切地说，它同构于 \(\operatorname {R}^1f_*(\mathcal {O}_X)\) 的 \(i\) 次外幂。\end{remark}\begin{remark}{3.8}\setnumber{3.8}\label{fga3.vi-3-remark-3.8}\skipsinglepar\par{}对于光滑射影态射 \(f\colon  X\to  S\)，若 \(S\) 约化且所有剩余特征为零，则 \hyperref[fga3.vi-3-corollary-3.6]{推论~\ref*{fga3.vi-3-corollary-3.6}} 的结果早已通过超越方法作为 Hodge 理论的推论而知道。事实上，此时所有 \(\operatorname {R}^pf_*(\Omega _{X/S}^q)\) 都局部自由。不过，在不等特征情形中，Serre 簇（[\hyperref[Ser1958b]{8}，第 50 页]）给出“\(\operatorname {R}^1f_*(\mathcal {O}_X)\) 局部自由”的反例。等特征情形中似乎尚无反例。\end{remark}\begin{corollary}{3.8}\setnumber{3.8}\label{fga3.vi-3-corollary-3.8}\skipsinglepar\par{}在 \hyperref[fga3.vi-3-theorem-3.5]{定理~\ref*{fga3.vi-3-theorem-3.5}} 的条件下，\(\underline {\operatorname {Pic}}_{X/S}^\sigma |U\) 在 \(\underline {\operatorname {Pic}}_{X/S}^\sigma |U\) 的所有点处在 \(U\) 上光滑。\begin{proof}\par{}应用 \hyperref[fga3.vi-2-proposition-2.10]{命题~\ref*{fga3.vi-2-proposition-2.10}} 的 (i)。\end{proof}\end{corollary}\par{}特别地，考虑到 \hyperref[fga3.vi-2-proposition-2.10]{命题~\ref*{fga3.vi-2-proposition-2.10}} 的 (iii)，有：\begin{corollary}{3.9}\setnumber{3.9}\label{fga3.vi-3-corollary-3.9}\skipsinglepar\par{}假设 \(S\) 的所有剩余特征为零。那么，\(\underline {\operatorname {Pic}}_{X/S}^\tau \) 在 \(S\) 上光滑。\end{corollary}\begin{remark}{3.10}\setnumber{3.10}\label{fga3.vi-3-remark-3.10}\skipsinglepar\par{}例如，由此推出，若 \(\underline {\operatorname {Pic}}_{X/S}\) 还在 \(S\) 上 \emph{固有}，则 \(f\) 的几何纤维的 Néron–Severi 挠群在 \(S\) 的每个连通分支上恒定（当 \(f\) 光滑且射影时，这一点用超越方法也很显然）。注意，直接利用 \hyperref[fga3.vi-2-theorem-2.5]{定理~\ref*{fga3.vi-2-theorem-2.5}}，可以更一般地证明：若 \(\underline {\operatorname {Pic}}_{X/S}\) 在 \(S\) 上固有（例如 \(f\colon  X\to  S\) 光滑且射影），且 \(q\) 是一个不同于 \(S\) 的各剩余特征的素数，则 \(X/S\) 的几何纤维的 Néron–Severi 挠群的 \(q\)-准素分量在 \(S\) 的每个连通分支上恒定。不过，在特征 \(p>0\) 时，对于挠群的 \(p\)-准素分量，情形就不再如此。然而，\(\underline {\operatorname {Pic}}_{X_s/k(s)}/\underline {\operatorname {Pic}}_{X_s/k(s)}^{00}=T_{X_s/k(s)}\) 在各 \(k(s)\) 上的秩仍有可能局部常值。当 \(S\) 约化时，可以证明，这等价于 \(\underline {\operatorname {Pic}}_{X/S}^{00}\) 存在，且 \(\underline {\operatorname {Pic}}_{X/S}\) 在 \(S\) 上 \emph{平坦}；并且只须对 \(S\) 是离散赋值环的谱的情形验证。这一点在我考察过的少数例子中确实成立；但由于 \(S\) 为 Artin 环的谱时，相应陈述错误（参见 \hyperref[fga3.vi-2-remark-2.9]{注~\ref*{fga3.vi-2-remark-2.9}} 和 \hyperref[fga3.vi-3-remarks-3.4]{注~\ref*{fga3.vi-3-remarks-3.4}}），我们也不能过于乐观。\end{remark}\chapter{Picard 概形的有限性定理}\thispagestyle{fancy}\label{fga3.vi-4}\par{}设 \(f\colon  X\to  S\) 是射影平坦态射，且 \(\underline {\operatorname {Pic}}_{X/S}\) 存在。“Hilbert 多项式”\(Q\in \mathbb {Q}[t]\) 使我们可以将 \(\underline {\operatorname {Pic}}_{X/S}\) 分解为开集 \(\underline {\operatorname {Pic}}_{X/S}^Q\) 的和。若不加进一步假设，例如保证 \(\underline {\operatorname {Pic}}_{X/S}\) 在 \(S\) 上分离的假设，则这些开集在 \(S\) 上一般并非有限型；当 \(X\) 是“一条退化为两条直线的圆锥曲线”时就得到反例。不过，若 \(f\) 是 \emph{可分} 的，且具有 \emph{不可约几何纤维}，则这可能成立。这个问题与 \(\underline {\operatorname {Pic}}_{X/S}^\tau \) 是否在 \(S\) 上有限型有关；后者可能不对 \(X/S\) 的纤维作任何假设就成立。当 \(f\colon  X\to  S\) 光滑时，注意到 \(\underline {\operatorname {Pic}}_{X/S}^\tau \) 包含在某个 \(\underline {\operatorname {Pic}}_{X/S}^Q\) 中（因为在非奇异射影簇上，对于除子，“挠”等价细于数值等价——实际上由 Matsusaka 的结果，二者相同）；因此，若各 \(\underline {\operatorname {Pic}}_{X/S}^Q\) 在 \(S\) 上有限型，它也有限型。注意，刚才提出的这些有限性问题，即使不假设 \(\underline {\operatorname {Pic}}_{X/S}\) 存在，仍有意义，因为可以将它们表述为某些可逆模族在 \hyperref[fga3.iv]{[3]，IV} 的意义下“有界”。对于每个代数闭域 \(k\)，考虑 \(\mathbb {P}_k^r\) 中维数为 \(n\)、次数为 \(d\) 的整子概形，以及这些预概形上 Hilbert 多项式为 \(Q\) 的可逆模（其中 \(r\)、\(n\)、\(d\)、\(Q\) 均给定）；要证明这些模构成的族（视为各 \(\mathbb {P}_k^r\) 上的凝聚模）是 \emph{有界} 的，即可由 \(\mathbb {Z}\) 上的有限型概形参数化。\par{}利用 Matsusaka 的方法 \hyperref[Mat1957]{[6]}，一个相当技术性的证明（以适当形式使用“等价判据”）给出：若只考虑 \(\mathbb {P}_k^r\) 的非奇异子簇，答案为肯定。更确切地说，有下述结果：\begin{theorem}{4.1}\setnumber{4.1}\label{fga3.vi-4-theorem-4.1}\skipsinglepar\par{}设 \(f\colon  X\to  S\) 是几何纤维连通的光滑射影态射，\(S\) Noether。设 \(\mathcal {O}_X(1)\) 是 \(X\) 上相对于 \(S\) 的甚丰沛模，\(E\) 是 \(\underline {\operatorname {Pic}}_{X/S}\) 的子集，对应于 \(X/S\) 的几何纤维 \(\overline {X}_{s_i}\) 上的可逆模族 \((\mathscr {L}_i)\)；设 \(D_i\) 是 \(\overline {X}_{s_i}\) 上定义 \(\mathscr {L}_i\) 的除子（未必为正），\(a_n^{(i)}t^n+\ldots +a_0^{(i)}\) 是 \(\mathscr {L}_i\) 的 Hilbert 多项式，\(\xi =\xi _i\) 是定义 \(\mathcal {O}_X(1)\) 的除子，即一个超平面截面。\par{}那么，下列条件等价：\begin{enumerate}
    
 \item[a.]\(Q\) 拟紧，即包含在 \(S\) 上某个有限型开集中，也即族 \((\mathscr {L}_i)\) 有界。

    
 \item[b.]\(E\) 包含在有限多个集合 \(\underline {\operatorname {Pic}}_{X/S}^Q\)（\(Q\in \mathbb {Q}[t]\)）之并中，即各 \(\mathscr {L}_i\) 的 Hilbert 多项式构成的集合有限。

    
 \item[b'.]（若 \(X/S\) 的所有纤维具有相同维数 \(n\)。）各 \(\mathscr {L}_i\) 的 Hilbert 多项式的系数 \(a_{n-1}^{(i)}\) 和 \(a_{n-2}^{(i)}\) 包含在某个有限集中。

    
 \item[b''.]（若 \(X/S\) 的所有纤维具有相同维数 \(n\)。）整数 \(\xi ^{n-1}D_i\) 和 \(\xi ^{n-2}D_i^2\)（即 \(D_i\) 和 \(D_i^2\) 的射影次数）有上界。

  \end{enumerate}\end{theorem}\begin{corollary}{4.2}\setnumber{4.2}\label{fga3.vi-4-corollary-4.2}\skipsinglepar\par{}设 \(f\colon  X\to  S\) 是几何纤维连通的光滑射影态射。那么，概形 \(\underline {\operatorname {Pic}}_{X/S}^Q\)（\(Q\in \mathbb {Q}[t]\)）和 \(\underline {\operatorname {Pic}}_{X/S}^\tau \) 都在 \(S\) 上射影。\begin{proof}\par{}由 \hyperref[fga3.vi-4-theorem-4.1]{定理~\ref*{fga3.vi-4-theorem-4.1}}，它们在 \(S\) 上有限型，因此可应用 \hyperref[fga3.vi-2-corollary-2.4]{推论~\ref*{fga3.vi-2-corollary-2.4}}。\end{proof}\end{corollary}\begin{remark}{}\label{fga3.vi-4-remark-comp}\skipsinglepar\par{}\emph{[编者]} 自本报告编写以来，本节所提出的这一类有限性问题几乎已全部解决。这里指出目前已知的主要事实。为简化陈述，以下默认为各陈述中出现的所有 Picard 预概形都存在，尽管对陈述作显然的修改，就可以去掉任何存在性假设。以下 \(S\) 表示 Noether 概形，\(X\) 和 \(Y\) 是 \(S\) 上的固有概形。\begin{enumerate}
      
 \item[i.]设 \(f\colon  X\to  Y\) 是满射态射。那么 \(f^*\colon \underline {\operatorname {Pic}}_{Y/S}\to \underline {\operatorname {Pic}}_{X/S}\) 是有限型态射。


      
 \item[ii.]假设 \(Y\) 在 \(S\) 上射影，配有相对于 \(S\) 的丰沛可逆模，设 \(X\) 是此模的任意一个截面的零点概形。令 \(f\colon  X\to  Y\) 为典范浸入。最后，假设 \(Y/S\) 的纤维的不可约分支的维数为 \(\geqslant 3\)。那么 \(f^*\colon \underline {\operatorname {Pic}}_{Y/S}\to \underline {\operatorname {Pic}}_{X/S}\) 为有限型。


      
 \item[iii.]假设 \(X\) 在 \(S\) 上射影，且其所有几何纤维整，维数为 \(n\)。设 \(\mathcal {O}_X(1)\) 是 \(X\) 上的丰沛可逆模，用它来定义 Hilbert 多项式。设 \(M\) 是 \(\underline {\operatorname {Pic}}_{X/S}\) 的子集。那么，\(M\) 拟紧，当且仅当 \(M\) 的各元素的 Hilbert 多项式 \(a_0x^n+a_1x^{n-1}+a_2x^{n-2}+\ldots +a_n\) 中，系数 \(a_1\) 和 \(a_2\) 有界。


      
 \item[iv.]对每个整数 \(n\neq 0\)，群预概形 \(\underline {\operatorname {Pic}}_{X/S}\) 中的 \(n\) 次幂同态是有限型态射。

    \end{enumerate}\par{}注意，(i) 和 (ii) 还蕴含（在各自陈述的假设下）\emph{\(\underline {\operatorname {Pic}}_{Y/S}\) 的子集 \(M\) 拟紧，当且仅当它在 \(\underline {\operatorname {Pic}}_{X/S}\) 中的像拟紧}。因此，\(Y\) 上的可逆模 \(\mathscr {L}\) 与 \(0\) 是 \(\tau \)-等价的，当且仅当它在 \(X\) 中的逆像亦然；换言之，\(\underline {\operatorname {Pic}}_{Y/S}^\tau \) 是 \(\underline {\operatorname {Pic}}_{X/S}^\tau \) 的逆像。特别地，要证明第一个预概形在 \(S\) 上有限型，只须证明第二个如此，因为 \(\underline {\operatorname {Pic}}_{Y/S}\to \underline {\operatorname {Pic}}_{X/S}\) 为有限型。再利用 (i)、Chow 引理及 (iii)，得到：\begin{enumerate}
      
 \item[v.]\(\underline {\operatorname {Pic}}_{X/S}^\tau \) 在 \(S\) 上有限型。

    \end{enumerate}\par{}一般地，将有限态射的情形中的 (i) 与 (ii) 合用，对于大部分有限性问题，可以归约到 \(X/S\) 的几何纤维整、正规且维数为 \(\leqslant 2\) 的情形；往往甚至可以将满射但未必有限的态射的情形中的 (i)，与代数曲面的解奇点（在任意特征下由 Abhyankar 给出）合用，归约到 \(X/S\) 甚至光滑、从而具有维数为 \(2\) 的非奇异几何纤维的情形。例如，结合 (v) 及界定非奇异射影曲面的 Néron–Severi 群的秩的 Picard–Igusa 不等式，可证明 Néron 有限性定理的下述推广：\begin{enumerate}
      
 \item[vi.]设 \(X/S\) 在 \(S\) 上固有，除此之外任意。那么，\(X/S\) 的几何纤维 \(X_i/k_i\) 的 Néron–Severi 群 \(\underline {\operatorname {Pic}}_{X_i/k_i}/\underline {\operatorname {Pic}}_{X_i/k_i}^0\) 都是有限生成的，且它们的秩以及挠子群的阶有界。

    \end{enumerate}\par{}同样，归约到非奇异曲面的情形，再利用此情形中的已知定理（如 Néron 有限性定理及 Néron–Severi 群上的相交型非退化这一事实），便得到：\begin{enumerate}
      
 \item[vii.]设 \(X\) 是代数闭域上的固有概形。那么，\(X\) 中存在有限条整闭曲线 \(C_i\)（\(1\leqslant  i\leqslant  r\)），满足下列性质：对 \(\underline {\operatorname {Pic}}_{X/k}\) 的每个子集 \(M\)，\(M\) 拟紧，当且仅当整数 \(\deg \mathscr {L}_{C'_i}\)（\(\mathscr {L}\in  M\)）有界（其中 \(C'_i\) 表示 \(C_i\) 的正规化）。

    \end{enumerate}在上述陈述中，可以取 \(r\) 为 Néron–Severi 群的秩。一旦知道此群为有限型，(vii) 就归结为下述事实：由 \(X\) 中的曲线 \(C\) 定义的 Néron–Severi 群上的线性型，只有在 Néron–Severi 群的挠元上才同时为零。在 \(X\) 非奇异且射影的情形中，这一结果与 (v) 都由 Matsusaka 给出。利用 (vii)，容易得到 \(X\) 上与 \(0\) 为 \(\tau \)-等价的可逆模的下述刻画：\begin{enumerate}
      
 \item[viii.]
        \par{}设 \(X/k\) 是一个域上的固有概形，\(\mathscr {L}\) 是 \(X\) 上的可逆模。那么下列条件等价：\begin{enumerate}
            
 \item[a.]\(\mathscr {L}\) 与 \(0\) 是 \(\tau \)-等价的。

            
 \item[b.]对每个凝聚模 \(F\)，有 \(\chi (F\otimes \mathscr {L})=\chi (F)\)，其中 \(\chi \) 表示 Euler–Poincaré 特征。

            
 \item[b'.]与 (b) 相同，但取 \(F=\mathcal {O}_Y\)，其中 \(Y\) 是 \(X\) 中维数为 \(1\) 的整闭子概形。

            
 \item[c.]对上述每个 \(Y\)，以 \(Y'\) 表示其正规化曲线，有 \(\deg \mathscr {L}_{Y'}=0\)。

          \end{enumerate}若 \(X/k\) 射影且配有丰沛可逆模 \(\mathcal {O}_X(1)\)，上述条件还等价于：\begin{enumerate}
            
 \item[d.]对每个整数 \(m\)，\(\mathscr {L}^{\otimes  m}(1)\) 都丰沛。

            
 \item[e.]（若 \(X\) 整。）对每对整数 \((m,n)\)，有 \[                 \chi (\mathscr {L}^{\otimes  m}(n))                 = \chi (\mathcal {O}(n))               \]（即 (b) 对 \(F=\mathscr {L}^{\otimes  m}(n)\) 成立）。

          \end{enumerate}
      

    \end{enumerate}\par{}为证明最后一个条件的充分性，注意它蕴含各 \(\mathscr {L}^{\otimes  m}\) 的 Hilbert 多项式全都相等；于是由 Mumford 判据 (iii)，各 \(\mathscr {L}^{\otimes  m}\) 留在 \(\underline {\operatorname {Pic}}_{X/k}\) 的某个拟紧子集中，即满足 (a)。条件 (b)、(b') 和 (c) 应视为通常对非奇异射影簇定义的 \emph{数值等价} 概念在任意固有概形上的变式。对于这样的簇，(a) 与 (c) 的等价性显然早已知道（Matsusaka）。\par{}上述判据 (e) 还蕴含下列结果：\begin{enumerate}
      
 \item[ix.]设 \(f\colon  X\to  S\) 是几何纤维整的平坦射影态射。那么 \(\underline {\operatorname {Pic}}_{X/S}^\tau \) 在 \(\underline {\operatorname {Pic}}_{X/S}\) 中开 \emph{且闭}。

    \end{enumerate}\par{}我们只就关键结果 (i)、(ii)、(iii) 的证明作一些评注（结果 (iv) 与其他结果稍有不同，仅用根式满射有限态射的情形中的 (i)，更确切地说，仅用 Frobenius 态射的情形，就可证明）。对于 (i)，非平坦下降的思想（见 \hyperref[fga3.i-a.2.c]{TDTE I，第9页}）起着本质作用。可以发现，（若只考虑有限性结果，）缺少下降数据的有效性判据并无妨碍。Mumford 最近证明了 (iii) 的一个稍弱形式，其中判据用到 Hilbert 多项式的 \emph{全部} 系数。他的论证极为简单，受到 Nakai 的一个丰沛性判据的证明的启发（Nakai 对非奇异曲面陈述了该判据，Mumford 将它推广到任意射影态射）。在我看来，这一论证需要对 \(X/S\) 的纤维附加一个温和的限制（Serre 的（\(\text {S}_2\)）性质）；当几何纤维正规时，此限制得到满足。然后，在 (ii) 的证明中使用这一受限判据：判据 (i) 使我们可以归约到 \(Y/S\) 平坦、几何纤维整且正规的情形；应用 Mumford 判据，又容易归约到 \(X/S\) 满足相同条件的情形。由维数假设，\(Y\) 和 \(X\) 的几何纤维在其闭点处的深度为 \(\geqslant 2\)；这就可以按 [\hyperref[fga3.vi-bib-5]{5}，XII] 中给出的形式应用“等价判据”，从而完成 (ii) 的证明。有了 (i) 和 (ii) 以后，在 Mumford 判据中去掉纤维正规的假设，并证明 (iii) 所给的更强形式，就不困难了。最后注意，(i) 和 (ii) 的证明还表明，当 \(S\) 是域 \(k\) 的谱时，态射 \(\underline {\operatorname {Pic}}_{Y/k}\to \underline {\operatorname {Pic}}_{X/k}\) 是 \emph{仿射} 的（不仅仅是有限型的）。\end{remark}
\section*{原报告参考文献}
\begin{enumerate}
\item[{[1]}]\setnumber{1}\label{fga3.vi-bib-1} Chevalley (Claude). \emph{Sur la théorie de Picard}. Amer. J. of Math., t.~82, 1960, pp.~435--490.
\item[{[2]}]\setnumber{2}\label{fga3.vi-bib-2} Grothendieck (A.) et Dieudonné (J.). \emph{Éléments de géométrie algébrique. Chapitre I et suivants}. Paris, Presses universitaires de France, 1960--1961 (Institut des Hautes Études Scientifiques, Publications mathématiques 4, 8, 11, \ldots).
\item[{[3]}]\setnumber{3}\label{fga3.vi-bib-3} Grothendieck (Alexander). \emph{Technique de descente et théorèmes d'existence en géométrie algébrique, I--V}. Séminaire Bourbaki, t.~12, 1959/60, n\textsuperscript{o}~190, 29~页；n\textsuperscript{o}~195, 22~页；t.~13, 1960/61, n\textsuperscript{o}~212, 20~页；n\textsuperscript{o}~221, 28~页；t.~14, 1961/62, n\textsuperscript{o}~232, 19~页。
\item[{[4]}]\setnumber{4}\label{fga3.vi-bib-4} Grothendieck (Alexander). \emph{Séminaire de géométrie algébrique, 1960/61}. Paris, Institut des Hautes Études Scientifiques, 1961.
\item[{[5]}]\setnumber{5}\label{fga3.vi-bib-5} Grothendieck (Alexander). \emph{Séminaire de géométrie algébrique, 1961/62}，由一组听众撰写。Paris, Institut des Hautes Études Scientifiques（待出版）。
\item[{[6]}]\setnumber{6}\label{fga3.vi-bib-6} Matsusaka (T.). \emph{The criteria for algebraic equivalence and the torsion group}. Amer. J. Math., t.~79, 1957, pp.~53--66.
\item[{[7]}]\setnumber{7}\label{fga3.vi-bib-7} Mumford (D.) and Tate (J.). \emph{Séminaire de géométrie algébrique}. Harvard University, Spring Term 1962（待出版）。
\item[{[8]}]\setnumber{8}\label{fga3.vi-bib-8} Serre (Jean-Pierre). \emph{Sur la topologie des variétés algébriques en caractéristique $p$}. \emph{Symposium internacional de topologia algebraica} [1956, Mexico], pp.~24--53. Mexico, Universidad nacional autonoma, 1958.
\end{enumerate}
